\documentclass[10pt,letterpaper]{article}
\usepackage[T1]{fontenc}
\usepackage{amsmath,amsthm}
\usepackage{newpxtext,newpxmath}
\usepackage[margin=.88in]{geometry}
\usepackage{microtype,xcolor}
\definecolor{olive}{HTML}{505A3C}
\usepackage[colorlinks=true,urlcolor=olive,linkcolor=olive,
 pdftitle={Increasing Hahn series cannot be evaluated at omega}]{hyperref}
\setlength{\parindent}{0pt}
\setlength{\parskip}{.55em}
\pagestyle{empty}
\newcommand{\No}{\mathbf{No}}
\newcommand{\Q}{\mathbb Q}
\begin{document}
{\Large\bfseries\color{olive}Increasing Hahn series cannot be evaluated at $\omega$}\par
{\small A negative answer to Lipparini's Problem 7.7\hfill20 September 2026}
\vspace{.6em}

\textbf{The stated problem.}
Let $H$ be the ring of formal Hahn series
\[
 \sum_{\beta<\alpha}r_\beta x^{a_\beta},\qquad
 r_\beta\in\mathbb R\setminus\{0\},\quad
 \gamma<\beta\Longrightarrow a_\gamma<a_\beta,
\]
with ordinal-length increasing surreal exponent supports, coefficientwise
addition, and Hahn convolution. Problem 7.7 of [1] asks whether a ring
homomorphism $\Phi:H\to\No$ can satisfy $\Phi(x^a)=\omega^a$ for every
surreal $a$, with Conway's omega-map on the right.

\textbf{Theorem.} \emph{No such homomorphism exists. In fact, no unital
homomorphism from $\Q[[x]]$ into an ordered field can send $x$ to an element
greater than $1$.}

\textit{Proof.}
The subring $\Q[[x]]$ is contained in $H$, since every subset of
$\{0,1,2,\ldots\}$ is well-ordered in the required increasing direction.
Define $B(x)=\sum_{n\ge0}b_nx^n$ by
\[
 b_0=1,\qquad b_1=-\tfrac12,\qquad
 b_n=-\tfrac12\sum_{k=1}^{n-1}b_kb_{n-k}\quad(n\ge2).
\]
This is a well-defined formal series with rational coefficients. The constant
and linear coefficients of $B^2$ are $1$ and $-1$. At every degree $n\ge2$,
\[
 [x^n]B^2=2b_n+\sum_{k=1}^{n-1}b_kb_{n-k}=0.
\]
Consequently $B^2=1-x$. If $u$ is the image of $x$ under a unital homomorphism,
then
\[
 \Phi(B)^2=\Phi(B^2)=\Phi(1-x)=1-u.
\]
For $u>1$ this is a negative square, impossible in an ordered field. In the
stated problem $u=\omega>1$, and unitality follows from the prescribed image
of $x^0$.\hfill$\square$

\textbf{Why this addresses ordinary ring homomorphisms.}
No step asserts $\Phi(\sum b_nx^n)=\sum b_n\Phi(x)^n$.
The map is applied only to the finite algebraic identity $B^2+x-1=0$.
Continuity, order preservation, preservation of infinite sums, and fixation
of arbitrary real coefficients are not assumed. Proper-class issues are
irrelevant to this witness.

\textbf{Stronger consequence.}
For every positive integer $n$, the series $1\pm nx$ is a square of a unit
in $\Q[[x]]$, by the same recursion or substitution in $B$.
Its image is strictly positive, so every unital homomorphism into an ordered
field satisfies
\[
 -\frac1n<\Phi(x)<\frac1n\qquad(n\ge1).
\]
Thus $x$ must map to zero or an infinitesimal, never to $\omega$.
The accompanying article proves the converse for surreal targets and the
classification of pure monomial prescriptions for general Hahn series.

\vfill
{\small
\textbf{[1]} Paolo Lipparini, \emph{Monotone infinitary operations on ordinals
(extended version)}, \href{https://arxiv.org/abs/2505.00424v2}{arXiv:2505.00424v2},
30 April 2026, Problem 7.7, p.~38.

The question is explicit in the inspected revision. A focused search on
20 September 2026 did not locate a prior resolution; priority is not certified.
This independent argument has not been refereed or proof-assistant checked.}
\end{document}
